\section{Lecture 16 (November 6th)}
\begin{rmk}
Last class we have learned normal extensions. This class, we learn more about normal extensions and newly finite fields.
\end{rmk}
\vspace{2ex}
\begin{thm}
Let $F/K$ be an algebraic extension and $\bar{K}$ be an algebraic closure of $K$ with $F\subset \bar{K}$. The following are equivalent:
\begin{itemize}
\item[(i)] There exists $S\subset K[t]$ such that $F$ is a splitting field for $S$ (that is, $F=K(S_0)$ where $S_0$ is the set of zeros of $f(t)\in S$).
\item[(ii)] For each $K$-embedding $\sigma :F\rightarrow \bar{K}$, $\sigma (F)=F$.
\item[(iii)] $F/K$ is normal. That is, for each $\alpha \in F$, all zeros of $\mathrm{irr}(\alpha ,K)$ in $\bar{K}$ lie in $F$. 
\end{itemize}
\end{thm}
\vspace{2ex}
\begin{ex}
(14.13) Consider $p(t)=t^{4}+t^{3}+t^2+t+1\in {\bm Q}[t]$. Let $E={\bm Q}(\xi )$, $\xi $ is a primative 5th root of unity. Observe that $\sigma _{i}$ can be defined to send any $\xi \mapsto \xi ^{i}$ for $1\leq i\leq 4$.  
\begin{center}
\begin{tikzpicture}[baseline=(current bounding box.center),
  >={Straight Barb[length=3pt,width=6pt]}, % barbed arrowhead (side lines)
  node font=\small]

  % corners
  \node (F) at (3,4.5) {$\bar{{\bm Q}}$};
  \node (A) at (0,3) {$E={\bm Q}(\xi )$};
  \node (B) at (0,0) {${\bm Q}$};
  \node (C) at (3,0) {${\bm Q}$};
  \node (D) at (3,3) {${\bm Q}(\xi ^{i})={\bm Q}(\xi )$};

  % vertical arrows (pointing up)
  \draw[->] (B) -- (A);
  \draw[->] (C) -- (D);
  \draw[->] (D) -- (F);
  \draw[->] (A) -- node[midway,above] {$\sigma _{i}$}(D);

  % horizontal arrows (pointing right) with labels above
  \draw[->] (B) --  (C);

\end{tikzpicture}
\end{center}

Now consider $p(t)=t^3-2\in {\bm Q}[t]$. Let $E={\bm Q}(2^{1/3},\xi )$ where $\xi $ is a primative 3rd root of unity.

\begin{center}
\begin{tikzpicture}[baseline=(current bounding box.center),
  >={Straight Barb[length=3pt,width=6pt]}, % barbed arrowhead (side lines)
  node font=\small]

  % corners
  \node (A) at (0,2.5) {$F={\bm Q}(2^{1/3})$};
  \node (B) at (0,0) {${\bm Q}$};
  \node (C) at (3,0) {${\bm Q}$};
  \node (D) at (3,2.5) {$\bar{{\bm Q}}$};

  % vertical arrows (pointing up)
  \draw[->] (B) -- (A);
  \draw[->] (C) -- (D);
  \draw[->] (A) -- node[midway,above] {$\sigma $}(D);

  % horizontal arrows (pointing right) with labels above
  \draw[->] (B) --  (C);

\end{tikzpicture}
\end{center}
Here, $\sigma (2^{1/3})=2^{1/3}\xi \notin {\bm Q}(2^{1/3})$. 
\end{ex}
\vspace{2ex}
\begin{cor}
Let $F/K$ be a finite extension. Then $F$ is normal over $K$ if and only if
\[[F:K]_{s}=|\mathrm{Aut}(F/K)|\]
where $\mathrm{Aut}(F/K)$ is the group of all $K$-isomorphisms $\sigma :F\rightarrow F$, $\sigma \cdot \mu =\sigma \circ \mu $. 
\end{cor}
\vspace{2ex}
\begin{rmk}
(Observation) Choose an algebraic closure of $K$ with $K\subset F\subset \bar{K}$. There is an injection
\[\mathrm{Aut}(F/K)\rightarrow \{\sigma :F\rightarrow \bar{K}\}\]
where $\sigma $ is a $K$-embedding. Schematically, we have a map 
\newline
\begin{center}
$\Bigg($
\begin{tikzpicture}[baseline=(current bounding box.center),
  >={Straight Barb[length=3pt,width=6pt]}, % barbed arrowhead (side lines)
  node font=\small]

  \node (A) at (1.5,2) {$F$};
  \node (B) at (-1.5,2) {$F$};
  \node (C) at (0,0) {$K$};

  \draw[->] (B) --node[midway,above] {$\sigma $} (A);
  \draw (A) -- (C);
  \draw (C) -- (B);

\end{tikzpicture}
$\Bigg)$
$\longmapsto$
$\Bigg($
\begin{tikzpicture}[baseline=(current bounding box.center),
  >={Straight Barb[length=3pt,width=6pt]}, % barbed arrowhead (side lines)
  node font=\small]

  \node (A) at (1.5,2) {$F$};
  \node (B) at (0,2) {$F$};
  \node (D) at (3,2) {$\bar{K}$};
  \node (C) at (1.5,0) {$K$};

  \draw[->] (B) --node[midway,above] {$\sigma $} (A);
  \draw (A) -- (C);
  \draw (C) -- (B);
  \draw[->] (A) --node[midway,above] {$\iota$} (D);
  \draw (C) -- (D);

\end{tikzpicture}
$\Bigg)$
\end{center}
\end{rmk}
\vspace{2ex}
\begin{prop}
If $F/K$ is a field extension of degree 2, then $F/K$ is normal over $K$. 
\end{prop}
\vspace{2ex}
\begin{proof}
Choose an element $\alpha \in F\backslash K$. Then $F=K(\alpha )$. Let $\mathrm{irr}(\alpha ,K)=t^2+at+b\in K[t]$. Note that
\[t^2+at+b=(t-\alpha )(t+(a+\alpha ))\in \bar{K}[t]\]
We claim that $F$ is a splitting field for $\mathrm{irr}(\alpha ,K)$. This follows from the observation that $-(a+\alpha )$ and $\alpha $ belong in $F$. An observation is that for $F/E/K$, if $F/K$ is normal, then $F/E$ is normal.
\end{proof}
\vspace{2ex}
\section{Chapter 15. Galois Theory}
\begin{defi}
(15.1) Let $F/K$ be a field extension. Then the Galois group of $F/K$ is the group of $K$-automorphisms on $F$. We will denote it by $\mathrm{Gal}(F/K)$.
\end{defi}
\vspace{2ex}

